% ============================================================================
% (A) TABELA 2.1 com limites de célula visíveis
%     Substituir a tabela actual por esta. Requer \usepackage{array} no preâmbulo
%     (quase de certeza já lá está; se não, acrescentar).
% ============================================================================

\begin{table}[htbp]
\centering
\caption{Alternativas de leitura do número consideradas e razão do seu descarte.}
\label{tab:alternativas}
\renewcommand{\arraystretch}{1.35}
\setlength{\tabcolsep}{6pt}
\begin{tabular}{|>{\raggedright\arraybackslash}p{3.6cm}|p{9.0cm}|}
\hline
\textbf{Abordagem} & \textbf{Razão do descarte} \\
\hline
EasyOCR e OCR clássico &
Mesma limitação estrutural do MMOCR, falha em números pequenos ou inclinados e não
distingue o número de outros elementos gráficos. \\
\hline
PaddleOCR (com \emph{fine-tuning}) &
Falha de segmentação (\emph{segfault}) ao inicializar no mesmo processo que o CUDA do
\emph{pipeline}, e mesmo isolado não recuperou leituras face ao Qwen. \\
\hline
HunyuanOCR &
\emph{VLM} mais pesado, sem ganho de exatidão que compensasse o custo de VRAM e de
latência no \emph{hardware} disponível. \\
\hline
Qwen2.5-VL-7B &
Requer cerca de 10\,GB de VRAM, acima dos 6\,GB da GPU, o que inviabiliza a sua
utilização. \\
\hline
Super-resolução (RealESRGAN) &
A informação em falta não existe na imagem de origem, e a ampliação amplifica o ruído
em vez de a recuperar. \\
\hline
SAM-3D-Body &
Reconstrução tridimensional do corpo demasiado pesada para correr por \emph{frame} e
desadequada a um simples sinal de orientação. \\
\hline
Classificador de orientação &
O ganho de filtrar as \emph{frames} sem as costas visíveis não compensou, pois a votação
por \emph{tracklet} já absorve naturalmente essas \emph{frames}. \\
\hline
\end{tabular}
\end{table}


% ============================================================================
% (B) INTRODUÇÃO E CONCLUSÃO COMO CAPÍTULOS NUMERADOS
% ============================================================================
%
% No ficheiro onde estão os \include/\input dos capítulos (normalmente main.tex),
% procurar as linhas da introdução e da conclusão. Vão ter uma destas formas:
%
%     \chapter*{INTRODUÇÃO}
%     \addcontentsline{toc}{chapter}{INTRODUÇÃO}
%
% ou uma macro do template, do género \introduction{} / \conclusions{}.
%
% CORREÇÃO, no primeiro caso:
%     \chapter{INTRODUÇÃO}                 % sem o asterisco
%     % apagar a linha \addcontentsline, que passa a ser automática
%
% CORREÇÃO, no segundo caso: abrir o .cls ou .sty do template, localizar a
% definição da macro e trocar lá o \chapter* por \chapter.
%
% A conclusão precisa ainda de duas correcções que já vinham do template:
%   1. o título está em inglês ("CONCLUSIONS"), deve passar a "CONCLUSÃO";
%   2. o corpo é ainda texto de exemplo em latim.
%
% ATENÇÃO À RENUMERAÇÃO. Com a introdução a passar a capítulo 1, tudo desloca:
%     Cap. 1  INTRODUÇÃO
%     Cap. 2  ESTADO DA ARTE
%     Cap. 3  METODOLOGIA
%     Cap. 4  RESULTADOS
%     Cap. 5  CONCLUSÃO
%
% Há referências no texto que ficam erradas e têm de ser actualizadas:
%   - "ESTRUTURA DO DOCUMENTO", na introdução (texto novo mais abaixo);
%   - "seguindo a lógica exposta na Secção 1.4"  -> Secção 2.4
%   - "a recalibração descrita na Secção 2.3.2"  -> Secção 3.3.2
%   - "a arquitectura da baseline (Figura 1.2)"  -> Figura 2.2
%   - "a conclusão do Capítulo 3"                -> Capítulo 4
%   - "o Capítulo 3 reporta e discute"           -> Capítulo 4
%
% Para não voltar a passar por isto, vale a pena pôr \label{cap:...} em cada
% capítulo e substituir os números por \ref{}.


% ============================================================================
% (C) ESTRUTURA DO DOCUMENTO reescrita para a numeração nova
%     (secção final da introdução)
% ============================================================================

\section*{ESTRUTURA DO DOCUMENTO}

A dissertação está organizada em cinco capítulos. O presente capítulo enquadra o problema
e delimita o âmbito do trabalho. O Capítulo~\ref{cap:estado_arte} revê o estado da arte,
percorrendo primeiro a literatura sobre reconhecimento do número de camisola e atribuição
de equipa, e detendo-se depois sobre os sistemas apresentados nas duas edições do desafio
SoccerNet, que constituem o termo de comparação mais direto para este trabalho. O
Capítulo~\ref{cap:metodologia} descreve a metodologia proposta, detalhando o novo módulo
de leitura do número, as correções introduzidas na atribuição de equipa e do lado do campo
e o ambiente experimental em que tudo assenta. O Capítulo~\ref{cap:resultados} reporta e
discute os resultados, primeiro de cada módulo isoladamente e depois no sistema completo,
incluindo as experiências controladas que delimitam o teto da abordagem. O
Capítulo~\ref{cap:conclusao} sintetiza os contributos e aponta as direções de trabalho
futuro.
